\section{VR Sport and Rowing-Related Work}
\label{sec:background_vr_sport_rowing_related_work}

VR sport systems combine bodily action with simulated environments. They can be used for training, exercise, rehabilitation, motivation, or entertainment. Compared with passive VR experiences, sport and exercise systems place greater demands on the relation between movement, feedback, and physical effort. The user is not only observing a scene, but acting continuously within it.
% TODO cite: interactive VR sport systematic review.
% TODO cite: VR sport/training review.

A recurring issue in VR sport is fidelity. Fidelity can refer to the visual realism of the environment, the realism of the task, the correspondence between real and virtual movement, or the usefulness of the simulation for training transfer. These forms of fidelity are related, but they are not the same. A system can look realistic while simplifying movement, or reproduce movement constraints while presenting a more abstract environment.
% TODO cite: Richlan et al. or related VR sport fidelity/transfer source.
% TODO cite: Harris et al. on VR action/perception caution.

Rowing has appeared in several virtual or interactive training contexts. Some work has investigated VR or virtual environments for rowing exercise, training, social presence, or transfer between virtual and real rowing. These studies are useful as related work, but the present thesis addresses a more specific embodiment question. It does not primarily evaluate training effectiveness, physiological performance, or long-term motivation. Instead, it asks how two representation strategies affect the experienced relation between the user's body, movement, and virtual rowing action.
% TODO cite: VR rowing HMD workout source.
% TODO cite: presence of others during VR rowing exercise.
% TODO cite: virtual-to-real rowing skill learning / rowing simulator work if used.

Adjacent exercise systems such as VR cycling are also relevant because they show how stationary physical effort can be connected to virtual locomotion. In such systems, visual motion, environmental context, avatars, and interaction feedback can influence the exercise experience. However, rowing differs from cycling because the hands, handle, seat, trunk, and legs form a highly visible coordinated movement. The user's upper body and tool interaction are therefore more central to the embodied visual representation.
% TODO cite: VR cycling/exercise adjacent studies only if needed.

The contribution of this thesis is therefore not simply another VR rowing prototype. Its focus is the comparison between an ergometer-centred, high-congruence representation and a boat-centred, sport-depiction representation within the same general VR rowing context. This comparison separates two design values that are often blended together: close correspondence with the user's actual movement and convincing depiction of the sport scenario.

Related work in VR sport provides the broader design context, while embodiment research provides the conceptual lens. Rowing biomechanics explains why the two mappings are not equivalent. Together, these perspectives motivate the empirical question: when the user rows on an ergometer in VR, is embodiment better supported by preserving the felt movement geometry or by presenting a more recognizable on-water rowing action?
